\documentclass[../../../include/open-logic-section]{subfiles}

\begin{document}

\olfileid{sth}{spine}{recursion}
\olsection{అతిపరిమిత పునరావృత్తి సిద్ధాంతం (దాని రూపాలు)}

అయితే ముందుగా
\olref[valpha]{defValphas}
విజయవంతమైన నిర్వచనమని
నిర్ధారించాలి. మరింత
సాధారణంగా, అతిపరిమిత
పునరావృత్తి ద్వారా
నిర్వచనం ఇచ్చే ఏ
ప్రయత్నమైనా ఫలిస్తుందని
నిరూపించాలి. ఈ
విభాగం లక్ష్యం అదే.

\emph{హెచ్చరిక: ఇది
కష్టమైన విషయం}.
అయితే ప్రధాన సందేశం
చాలా సరళం: అతిపరిమిత
ఆగమనం, ప్రతిస్థాపన
కలిసి అతిపరిమిత
పునరావృత్తి (పలు
రూపాలు) చెల్లుబాటును
నిర్ధారిస్తాయి.\footnote{గుర్తుంచుకోండి:
వేరుగా స్పష్టంగా
చెప్పనంతవరకు, అన్ని
సూత్రాలు, పదాల్లో
పరామితులు ఉండవచ్చు.}

\begin{defn}
	$\tau(x)$ ఒక పదం,
	$f$ ఒక ప్రమేయం,
	$\alpha$ ఒక క్రమసంఖ్య
	అనుకుందాం.
	$\dom{f} = \alpha$
	మరియు
	$(\forall \beta \in \alpha)f(\beta) = \tau(\funrestrictionto{f}{\beta})$
	రెండూ అయితేనే
	$f$ను $\tau$కు
	\emph{$\alpha$-ఉజ్జాయింపు}
	అంటాం.
\end{defn}

\begin{lem}[పరిమిత పునరావృత్తి]\ollabel{transrecursionfun}
ఏ పదం $\tau(x)$,
ఏ క్రమసంఖ్య $\alpha$కైనా,
$\tau$కు ఏకైక
$\alpha$-ఉజ్జాయింపు
ఉంటుంది.
\end{lem}
\begin{proof}
ఏ $\gamma \leq \alpha$కైనా
ఏకైక $\gamma$-ఉజ్జాయింపు
ఉంటుందని చూపుతాం.

మొదట ఏకైకత్వాన్ని
స్థాపిద్దాం. $g$, $h$
వరుసగా $\gamma$-,
$\delta$-ఉజ్జాయింపులు
అనుకుందాం. వాటి
ఆర్గ్యుమెంట్లపై అతిపరిమిత
ఆగమనం ద్వారా,
ఏ $\beta \in \dom{g} \cap \dom{h} =
\gamma \cap \delta = \min(\gamma, \delta)$కైనా
$g(\beta) = h(\beta)$.
అందువల్ల మన ఉజ్జాయింపులు
(ఉన్నట్లయితే) ఏకైకాలు;
అన్ని సామాన్య విలువలపై
ఒకేలా ఉంటాయి.

ఇప్పుడు ఉనికిని
స్థాపించడానికి
$\delta \leq \alpha$
క్రమసంఖ్యలపై సరళ
అతిపరిమిత ఆగమనాన్ని
(\olref[ordinals][opps]{simpletransrecursion})
వాడతాం.

శూన్య ప్రమేయం స్పష్టంగా
$\emptyset$-ఉజ్జాయింపు.

$g$ ఒక
$\gamma$-ఉజ్జాయింపు
అయితే,
$g \cup
\{\tuple{\gamma, \tau(g)}\}$
ఒక $\ordsucc{\gamma}$-ఉజ్జాయింపు.

$\gamma$ సీమా క్రమసంఖ్య,
ప్రతి $\delta < \gamma$కు
$g_\delta$ ఒక
$\delta$-ఉజ్జాయింపు
అయితే, $g = \bigcup_{\delta \in \gamma} g_\delta$
అనుకుందాం. వివిధ
$g_\delta$లు అన్ని
సామాన్య విలువలపై
ఒకేలా ఉంటాయి
కాబట్టి ఇది ప్రమేయం.
ఇంకా $\delta \in \gamma$
అయితే $g(\delta) = g_{\ordsucc{\delta}}(\delta) =
\tau(\funrestrictionto{g_{\ordsucc{\delta}}}{\delta}) =
\tau(\funrestrictionto{g}{\delta})$.

దీనితో అతిపరిమిత
ఆగమన నిరూపణ పూర్తయింది.
\end{proof}

ప్రమేయానికి బదులు
\emph{పదాన్ని} నిర్వచించడానికి
అనుమతిస్తే, ముందరి
ఫలితంలోని $\alpha$
హద్దును తొలగించవచ్చు.
తదుపరి ఫలితం ప్రకటన,
నిరూపణలో $\sigma$
పదం అయితే
$\funrestrictionto{\sigma}{\alpha} = \Setabs{\tuple{\beta,
\sigma(\beta)}}{\beta \in \alpha}$
అనుకుందాం.

\begin{thm}[సాధారణ పునరావృత్తి]\ollabel{transrecursionschema}
ఏ పదం $\tau(x)$కైనా,
ప్రతి క్రమసంఖ్య~$\alpha$కు
$\sigma(\alpha) = \tau(\funrestrictionto{\sigma}{\alpha})$
అయ్యే పదం $\sigma(x)$ను
స్పష్టంగా నిర్వచించవచ్చు.
\end{thm}

\begin{proof}
ప్రతి $\alpha$కూ
\olref{transrecursionfun}
ప్రకారం $\tau$కు
ఏకైక $\alpha$-ఉజ్జాయింపు
$f_\alpha$ ఉంటుంది.
$\sigma(\alpha)$ను
$f_{\ordsucc{\alpha}}(\alpha)$గా
నిర్వచిద్దాం. ఇప్పుడు:
	\begin{align*}
		\sigma(\alpha) &= 
		f_{\ordsucc{\alpha}}(\alpha) \\&= 
		\tau(\funrestrictionto{f_{\ordsucc{\alpha}}}{\alpha}) \\&= 
		\tau(\Setabs{\tuple{\beta, f_{\ordsucc{\alpha}}(\beta)}}{\beta \in \alpha}) \\&= 
		\tau(\Setabs{\tuple{\beta, f_{\ordsucc{\beta}}(\beta)}}{\beta \in \alpha})\\&=
		\tau(\funrestrictionto{\sigma}{\alpha})
	\end{align*}
ఎందుకంటే
\olref{transrecursionfun}లోని
విధంగా, అన్ని
$\beta < \alpha$లకు
$f_{\ordsucc{\beta}}(\beta) = f_{\ordsucc{\alpha}}(\beta)$.
\end{proof}
\noindent
\olref{transrecursionschema}
ఒక \emph{పథకం}
అని గమనించండి.
ముఖ్యంగా, $\sigma$
ఒక ప్రమేయాన్ని,
అంటే ఒక రకమైన
\emph{సమితిని},
నిర్వచిస్తుందని
ఆశించలేం. అలా
అయితే $\dom{\sigma}$
అన్ని క్రమసంఖ్యల
సమితి కావాలి; ఇది
బురాలి-ఫోర్టీ వైరుధ్యానికి
(\olref[ordinals][basic]{buraliforti})
విరుద్ధం.

అయినా \olref{transrecursionschema}
మన $V_\alpha$ల
నిర్వచనాన్ని సమర్థిస్తుందని
ఇంకా చూపాలి.
అది వెంటనే స్పష్టం
కాకపోవచ్చు; కానీ
అతిపరిమిత పునరావృత్తి
చివరి, సరళ రూపంతో
స్పష్టమవుతుంది.

\begin{thm}[సరళ పునరావృత్తి]\ollabel{simplerecursionschema}
ఏ పదాలు $\tau(x)$,
$\theta(x)$, ఏ సమితి
$A$కైనా, ఈ కింది
షరతులు గల పదం
$\sigma(x)$ను స్పష్టంగా
నిర్వచించవచ్చు:
\begin{align*}
	\sigma(\emptyset) &= A\\
	\sigma(\ordsucc{\alpha}) &= \tau(\sigma(\alpha)) &&
		\alpha\text{ ఏ క్రమసంఖ్య అయినా}\\
	\sigma(\alpha) &= \theta(\ran{\funrestrictionto{\sigma}{\alpha}})&&
	\alpha\text{ సీమా క్రమసంఖ్య అయినప్పుడు}
\end{align*}
\end{thm}

\begin{proof}
మొదట $\xi(x)$
పదాన్ని ఇలా
నిర్వచిద్దాం:
%t $\xi(x) = A$ అనేది $x$ క్రమసంఖ్య నిర్వచన సమితిగల ప్రమేయం కాకపోతే;
%లేకపోతే ఇలా:
\[
	\xi(x) = 
	\begin{cases}
			A &\text{ఒకవేళ $x$ నిర్వచన సమితి క్రమసంఖ్యగల ప్రమేయం కాకపోతే}\\
			&\hspace{1em}\text{లేదా నిర్వచన సమితి శూన్యమైతే; లేకపోతే:}\\
			\tau(x(\alpha)) & \text{ఒకవేళ $\dom{x} = \ordsucc{\alpha}$}\\
			\theta(\ran{x}) & \text{ఒకవేళ $\dom{x}$ సీమా క్రమసంఖ్య అయితే}
	\end{cases}
\]
\sourcecorrection{OLTESTSPINREC-001}{మూల విభజిత నిర్వచనంలో శూన్య నిర్వచన సమితిగల ప్రమేయానికి ఏ శాఖా వర్తించదు, అయినా తరువాతి నిరూపణ దాని విలువను వాడుతుంది. మొదటి శాఖలో ఆ సందర్భాన్ని చేర్చాం.}
\olref{transrecursionschema}
ప్రకారం ప్రతి క్రమసంఖ్య
$\alpha$కు
$\sigma(\alpha) = \xi(\funrestrictionto{\sigma}{\alpha})$
అయ్యే పదం $\sigma(x)$
ఉంటుంది; ఇంకా
$\funrestrictionto{\sigma}{\alpha}$,
నిర్వచన సమితి
$\alpha$ అయిన
ప్రమేయం. $\sigma$కు
కావలసిన ధర్మాలు
ఉన్నాయని సరళ అతిపరిమిత
ఆగమనంతో
(\olref[ordinals][opps]{simpletransrecursion})
చూపుతాం.

మొదట
$\sigma(\emptyset) = \xi(\emptyset) = A$.

తరువాత
$\sigma(\ordsucc{\alpha}) = \xi(\funrestrictionto{\sigma}{\ordsucc{\alpha}}) = \tau(\funrestrictionto{\sigma}{\ordsucc{\alpha}}(\alpha)) = \tau(\sigma(\alpha))$.

చివరిగా $\alpha$
సీమా క్రమసంఖ్య అయితే,
$\sigma(\alpha) = \xi(\funrestrictionto{\sigma}{\alpha}) = \theta(\ran{\funrestrictionto{\sigma}{\alpha}})$.
\end{proof}
\noindent
ఇప్పుడు \olref[valpha]{defValphas}ను
సమర్థించడానికి
$A
= \emptyset$, $\tau(x) = \Pow{x}$,
$\theta(x) = \bigcup x$
అని తీసుకుంటే చాలు.
ఎట్టకేలకు $V_\alpha$ల
నిర్వచనం సమర్థించబడింది!{}


\end{document}
